\documentclass[a4paper,10pt]{article}
\usepackage[utf8]{inputenc}
\usepackage[T1]{fontenc}
\usepackage[french]{babel}
\usepackage{geometry}
\geometry{left=1.5cm,right=1.5cm,top=1.2cm,bottom=1.2cm}
\usepackage{amsmath,amssymb}
\usepackage{enumitem}
\usepackage{array}
\usepackage{booktabs}
\usepackage{xcolor}
\usepackage{tcolorbox}
\usepackage{tikz}
\usepackage{multicol}

\definecolor{bleupro}{HTML}{0969da}
\definecolor{orange}{HTML}{d97706}
\definecolor{vertpro}{HTML}{2f855a}
\definecolor{rouge}{HTML}{c53030}
\definecolor{gris}{HTML}{718096}

\tcbuselibrary{skins,breakable}

\newtcolorbox{defbox}[1]{
  colback=blue!3,
  colframe=bleupro,
  title={\small\textbf{#1}},
  fonttitle=\bfseries,
  boxsep=3pt,
  left=5pt,right=5pt,top=3pt,bottom=3pt
}

\newtcolorbox{methbox}[1]{
  colback=green!3,
  colframe=vertpro,
  title={\small\textbf{#1}},
  fonttitle=\bfseries,
  boxsep=3pt,
  left=5pt,right=5pt,top=3pt,bottom=3pt
}

\newtcolorbox{attbox}{
  colback=red!3,
  colframe=rouge,
  title={\small\textbf{Attention !}},
  fonttitle=\bfseries,
  boxsep=3pt,
  left=5pt,right=5pt,top=3pt,bottom=3pt
}

\newtcolorbox{exbox}{
  colback=yellow!5,
  colframe=orange,
  boxsep=3pt,
  left=5pt,right=5pt,top=3pt,bottom=3pt
}

\setlength{\parindent}{0pt}
\setlength{\parskip}{4pt}
\pagestyle{empty}

\begin{document}

\begin{center}
{\Large\bfseries\color{bleupro} Fiche de révision — CCF Mathématiques}\\[4pt]
{\normalsize Terminale Bac Pro ERA-TMA / ICCER \quad | \quad Probabilités \& Polynômes degré 3}
\end{center}

\vspace{6pt}
\begin{multicols}{2}

% ========================================
{\large\bfseries\color{bleupro} PROBABILITÉS}
\vspace{4pt}

\begin{defbox}{Arbre de probabilités pondéré}
\begin{itemize}[nosep,leftmargin=12pt]
  \item Chaque branche porte une probabilité
  \item La somme des branches d'un même nœud $= 1$
  \item Probabilité d'un chemin $=$ produit des branches
\end{itemize}
\begin{center}
\begin{tikzpicture}[scale=0.7, every node/.style={font=\scriptsize}]
  \node[circle,fill=black,inner sep=1pt] (R) at (0,0) {};
  \node (A) at (2,0.8) {$A$};
  \node (Ab) at (2,-0.8) {$\overline{A}$};
  \draw (R)--(A) node[midway,above,sloped]{$P(A)$};
  \draw (R)--(Ab) node[midway,below,sloped]{$P(\overline{A})$};
  \node (B1) at (4.5,1.3) {$B$};
  \node (B1b) at (4.5,0.3) {$\overline{B}$};
  \draw (A)--(B1) node[midway,above,sloped]{$P_A(B)$};
  \draw (A)--(B1b) node[midway,below,sloped]{$P_A(\overline{B})$};
  \node (B2) at (4.5,-0.3) {$B$};
  \node (B2b) at (4.5,-1.3) {$\overline{B}$};
  \draw (Ab)--(B2) node[midway,above,sloped]{$P_{\overline{A}}(B)$};
  \draw (Ab)--(B2b) node[midway,below,sloped]{$P_{\overline{A}}(\overline{B})$};
\end{tikzpicture}
\end{center}
\end{defbox}

\begin{defbox}{Probabilité conditionnelle}
$$P_A(B) = \frac{P(A \cap B)}{P(A)} \quad \text{avec } P(A) \neq 0$$
$P_A(B)$ = probabilité de $B$ \textbf{sachant que} $A$ s'est réalisé.
\end{defbox}

\begin{defbox}{Formule des probabilités totales}
$$\boxed{P(B) = P(A) \times P_A(B) + P(\overline{A}) \times P_{\overline{A}}(B)}$$
On \textbf{additionne} les probabilités de tous les chemins menant à $B$.
\end{defbox}

\begin{defbox}{Indépendance}
$A$ et $B$ sont indépendants si et seulement si :
$$P(A \cap B) = P(A) \times P(B)$$
La réalisation de $A$ ne change pas la probabilité de $B$.
\end{defbox}

\begin{methbox}{Méthode — Exercice type probas}
\begin{enumerate}[nosep,leftmargin=12pt]
  \item Identifier les événements ($D$, $S$, $\overline{D}$...)
  \item Traduire l'énoncé en probabilités
  \item Construire l'arbre (vérifier : somme $= 1$)
  \item Calculer $P(\text{chemin})$ $=$ produit des branches
  \item Probabilités totales $=$ somme des chemins
\end{enumerate}
\end{methbox}

\begin{exbox}
\textbf{Exemple rapide :} $P(A) = 0{,}3$, $P_A(B) = 0{,}5$, $P_{\overline{A}}(B) = 0{,}1$.\\
$P(A \cap B) = 0{,}3 \times 0{,}5 = 0{,}15$\\
$P(\overline{A} \cap B) = 0{,}7 \times 0{,}1 = 0{,}07$\\
$P(B) = 0{,}15 + 0{,}07 = 0{,}22$
\end{exbox}

\begin{attbox}
\begin{itemize}[nosep,leftmargin=12pt]
  \item $P(\overline{A}) = 1 - P(A)$ (toujours !)
  \item Ne pas confondre $P(A \cap B)$ et $P(A \cup B)$
  \item Multiplier le long d'un chemin, additionner entre chemins
\end{itemize}
\end{attbox}

\columnbreak

% ========================================
{\large\bfseries\color{bleupro} POLYNÔMES DE DEGRÉ 3}
\vspace{4pt}

\begin{defbox}{Fonction polynôme de degré 3}
$$f(x) = ax^3 + bx^2 + cx + d \quad (a \neq 0)$$
\textbf{Dérivée :} $f'(x) = 3ax^2 + 2bx + c$
\end{defbox}

\begin{defbox}{Lien dérivée — variations}
\begin{center}
\renewcommand{\arraystretch}{1.3}
\begin{tabular}{|c|c|}
\hline
$f'(x) > 0$ & $f$ \textbf{croissante} $\nearrow$ \\
\hline
$f'(x) < 0$ & $f$ \textbf{décroissante} $\searrow$ \\
\hline
$f'(x) = 0$ et change de signe & \textbf{extremum} \\
\hline
\end{tabular}
\end{center}
\end{defbox}

\begin{methbox}{Méthode — Étude de fonction}
\begin{enumerate}[nosep,leftmargin=12pt]
  \item Calculer $f'(x)$
  \item Résoudre $f'(x) = 0$ (factoriser si possible)
  \item Tableau de signes : signe de chaque facteur
  \item Règle des signes $\to$ signe de $f'(x)$
  \item Tableau de variations ($\nearrow$ ou $\searrow$)
  \item Calculer $f$ aux extremums
\end{enumerate}
\end{methbox}

\begin{methbox}{Tableau de signes (règle des signes)}
$(x - a)$ est négatif si $x < a$, nul si $x = a$, positif si $x > a$.\\[4pt]
\textbf{Produit de deux facteurs :}\\
$(+) \times (+) = +$ \quad $(-) \times (-) = +$\\
$(+) \times (-) = -$ \quad $(-) \times (+) = -$
\end{methbox}

\begin{exbox}
\textbf{Exemple :} $V(x) = 4x^3 - 120x^2 + 900x$\\
$V'(x) = 12x^2 - 240x + 900 = 12(x-5)(x-15)$\\[4pt]
\renewcommand{\arraystretch}{1.3}
\begin{tabular}{|c|c|c|c|c|c|}
\hline
$x$ & $0$ & & $5$ & & $15$ \\
\hline
$(x-5)$ & & $-$ & $0$ & $+$ & \\
\hline
$(x-15)$ & & $-$ & & $-$ & $0$ \\
\hline
$V'(x)$ & & $+$ & $0$ & $-$ & \\
\hline
$V$ & $0$ & $\nearrow$ & $2000$ & $\searrow$ & $0$ \\
\hline
\end{tabular}
\end{exbox}

\begin{defbox}{Nombre de solutions de $f(x) = c$}
On trace la droite horizontale $y = c$ sur le tableau de variations :
\begin{itemize}[nosep,leftmargin=12pt]
  \item $c >$ max $\to$ 1 solution (ou 0 selon les cas)
  \item $c =$ max $\to$ 2 solutions
  \item min $< c <$ max $\to$ 3 solutions
  \item $c =$ min $\to$ 2 solutions
  \item $c <$ min $\to$ 1 solution
\end{itemize}
\end{defbox}

\begin{attbox}
\begin{itemize}[nosep,leftmargin=12pt]
  \item $f'(a) = 0$ ne suffit PAS pour un extremum ! (ex : $x^3$ en $0$)
  \item Il faut que $f'$ \textbf{change de signe}
  \item Utiliser la calculatrice pour tracer, puis vérifier par le calcul
\end{itemize}
\end{attbox}

\end{multicols}

\end{document}
